\section{Introduction}
\label{sec:intro}

A familiar shape recurs across the history of delegated emergency authority.
A legislature, faced with a circumstance it characterizes as exceptional,
grants an executive officer a power that was not previously available within the ordinary structure of the legal order.
The grant is framed as time-bounded.
The authority is to be used only while the exceptional circumstance persists,
and is to lapse, by its own terms, once the circumstance abates.
The grant is sometimes accompanied by a sunset provision,
sometimes by a periodic reauthorization requirement,
sometimes by reporting obligations to a legislative committee,
and sometimes by judicial review of the executive's exercise of the granted authority.
In each case, the structural intuition motivating the grant is that the authority is borrowed,
not given outright;
that it must return to the legislature when the exception ends;
and that the legal order will resume its ordinary shape thereafter.

The structural intuition is, with regularity that does not appear to be coincidental, falsified.
Article~48 of the Weimar Constitution authorized the Reich President to take ``the necessary measures''
to restore public security in cases where it was ``substantially disturbed or endangered''
\cite{TODO_weimar_const_art48}.
The article required notification to the Reichstag and permitted the Reichstag to demand the suspension of any measure taken under it.
In its first twelve years of operation, the authority was invoked approximately 250 times \cite{TODO_kershaw_hitler}.
By the time of the Reichstag Fire Decree in February 1933, the authority had become a routine instrument of governance
rather than the exceptional remedy the constitutional text described \cite{TODO_caldwell_popular_sovereignty}.
The transition from emergency power to instrument of dictatorship occurred along a continuous trajectory.
There was no single moment at which the authority's character changed.
There was, instead, a sequence of admissions, each of which appeared, at the time of its admission,
to be a defensible exercise of a power constitutionally available to its holder.

The Authorization for Use of Military Force enacted on September 18, 2001 \cite{TODO_aumf_2001}
permits the President to use ``all necessary and appropriate force'' against those determined to have
authored or aided the September 11 attacks.
The text contains no sunset.
It contains no geographic limit.
It is, twenty-five years later, the primary legal authority for armed action by United States forces
in countries that did not exist as targets of policy when the authorization was enacted
\cite{TODO_brennan_center_aumf,TODO_chesney_aumf}.
The authorization has been cited as legal basis in at least eight named theaters by 2024
\cite{TODO_jaffer_aumf}.
Periodic legislative efforts to repeal and replace the 2001 AUMF have not produced enacted legislation.
The pattern is not a failure of legislative will so much as a structural feature of the original grant:
the authorization, having been enacted without a return condition that would force its lapse,
has no built-in mechanism by which it terminates.
The political coalition required to enact a replacement is the same coalition required to enact a repeal,
and that coalition has not assembled in twenty-five years.

The Foreign Intelligence Surveillance Act's emergency authorization regimes,
including the seventy-two hour emergency authorization in 50 U.S.C. \S~1881a(c) (Section~702),
permit the Attorney General to authorize a surveillance program in advance of judicial approval
where the Attorney General determines that an emergency exists.
Declassified Foreign Intelligence Surveillance Court opinions and Privacy and Civil Liberties Oversight Board
reports document that the emergency authorization mechanism has operated at a boundary of judicial review
that the original statutory text did not contemplate \cite{TODO_pclob_702_report,TODO_fisc_702_opinions}.
The structural defect is the same: a grant of authority intended for time-bounded use that contains no
mechanism by which the underlying surveillance is rolled back when the emergency lapses.
The information collected during the emergency window remains in the holdings of the collecting agency.
The targets do not return to a prior state.

Article~17 of the General Data Protection Regulation grants data subjects a right to erasure of their personal data
in specified circumstances \cite{TODO_gdpr_art17}.
The Court of Justice of the European Union's enforcement of the article has produced erasure orders against
search-engine providers and platform operators that, in form, are time-bounded executive acts of administrative agencies.
In practice, an erasure order does not contemplate a path back to the prior state.
A search-engine index from which a result has been removed is not, in any practical sense, restored to its prior shape
by the lapse of the order \cite{TODO_kuner_gdpr_erasure}.
The order is structurally irreversible at the moment of its admission,
even though no statutory text characterizes it as such.

Section~230 of the Communications Decency Act, although not in its primary form a takedown regime,
operates in conjunction with the Digital Millennium Copyright Act's notice-and-takedown procedure
\cite{TODO_dmca_512}
and a constellation of platform self-regulatory practices to produce content removals that share
the structural shape of a time-bounded executive act \cite{TODO_keller_takedowns}.
The removed content, once removed, is not statutorily required to be restored upon any later event.
The takedown is, again, structurally irreversible at the moment of its admission.

These five regimes do not share doctrinal substance.
The constitutional emergency power of a Weimar-era Reich President bears no doctrinal relation to the
content-moderation practices of a contemporary American social-media platform,
and the surveillance authorities of FISA Section~702 bear no doctrinal relation to either.
They share, however, a structural feature that this Article identifies and names:
the original grant of authority was made without a typed rollback witness.
The authority could be admitted without exhibiting, at the moment of admission,
a constructible path back to the state of affairs that existed before the authority was exercised.

The Article's central claim is that this missing structural element is not a peripheral defect.
It is the explanation for the ratcheting failure mode that has, in each of these regimes and others,
transformed a time-bounded emergency authority into a permanent feature of the legal order.
The claim is not that a typed rollback witness, had it been required at the time these authorities were enacted,
would have prevented every adverse exercise of the authorities.
It is the narrower and more defensible claim that the absence of the typed rollback witness made the ratcheting
pattern structurally available, and that its presence would have made the same pattern structurally unavailable.

The grammar that supplies the typed rollback witness is the grammar of bounded executive action,
developed in a companion technical paper that constructs a Lean-attestable substrate for receipt-bounded
polities \cite{TODO_chio_parent_paper}.
The companion paper's contribution is the formal substrate.
This Article's contribution is the application of the substrate's grammar to the doctrinal pattern of
delegated emergency authority,
the demonstration that the absence of the grammar is the structural explanation for the ratcheting pattern,
and the exhibition of one implementation drawn from the substrate.

Part~\ref{sec:pattern} develops the historical pattern in detail and situates it within the Schmitt-Agamben
tradition on the state of exception.
Part~\ref{sec:cases} presents five case studies, with explicit acknowledgement that the case studies vary in
how cleanly they fit the framing.
Part~\ref{sec:grammar} states the formal grammar and explains the role of the typed rollback witness.
Part~\ref{sec:substrate} cites the Lean substrate without re-deriving it.
Part~\ref{sec:implementation} exhibits one implementation as a working example.
Part~\ref{sec:limits} acknowledges the limits of the formal approach,
particularly the irreducibly political content of the exception that the wrapper cannot capture.
Part~\ref{sec:conclusion} concludes.
